\documentclass[12pt]{article}
\usepackage[a4paper, tmargin=3cm, lmargin=2.5cm, rmargin=2.5cm, bmargin=3cm]{geometry}
\usepackage{hyperref}

\usepackage[brazilian]{babel} % Remover linha em caso de escrita em inglês
\usepackage[fixlanguage]{babelbib} % Remover linha em caso de escrita em inglês


\hypersetup{
    colorlinks=true,
    linkcolor=black,
    filecolor=magenta,
    urlcolor=blue,
    citecolor=black,
}

\usepackage{natbib} % to cite references by surname and year
\usepackage{graphicx}
\usepackage{fancyhdr}
\pagestyle{fancy} % Ativar o estilo fancy
\fancyhf{} % Limpar cabeçalhos e rodapés
\renewcommand{\headrulewidth}{0pt} % Remover a linha do cabeçalho
\fancyfoot[C]{IC/UNICAMP \\ 2026 }


%%%%%%%%%%%%%%%%%%%%%%%%%%%%%%%%%%
% CONFIGURAÇÃO DA PRIMEIRA PÁGINA
%%%%%%%%%%%%%%%%%%%%%%%%%%%%%%%%%%


\begin{document}

\vspace{5mm}
\hrule

\vspace{3mm}

\begin{center}
    {\large\textbf{INSTITUTO DE COMPUTAÇÃO - UNICAMP}} \\
    \vspace{3mm}
    {\large\textbf{MO810A/MC959A - Sistemas Multiagentes com Modelos de Linguagem}} \\
    \vspace{3mm}
    {\large\textbf{Prof. Dr. Julio Cesar dos Reis}}
\end{center}

\vspace{3mm}
\hrule

\vspace{2.5cm}
\begin{flushright}
    \Large{\textbf{Fase 2 - Metodologia}} \\
    \vspace{1cm}
    \textit{
        $<$Título do projeto$>$ \\
    }
    \vspace{5cm}
    $<$Aluno(a)$>$ - $<$RA Aluno(a)$>$ \\
    $<$Aluno(b)$>$ - $<$RA Aluno(b)$>$
\end{flushright}

%%%%%%%%%%%%%%%%%%%%%%%%%%%%%%%%%%
%%%%%%%%%%%%%%%%%%%%%%%%%%%%%%%%%%


\newpage
\pagestyle{fancy} % 2. Ativar o estilo fancy
\fancyhf{} % 3. Limpar cabeçalhos e rodapés
\renewcommand{\headrulewidth}{0pt} % Remover a linha do cabeçalho

\fancyfoot[R]{\thepage}


%%%%%%%%%%%%%%%%%%%%%%%%%%%%%%%%%%
%       INÍCIO DAS SEÇÕES
%%%%%%%%%%%%%%%%%%%%%%%%%%%%%%%%%%

%%%%%%%%%%%%%%%%%%%%%%%%%%%%%%%%%%%%%%%%
\section{Visão Geral da Metodologia}
\label{sec:visao_geral}
%%%%%%%%%%%%%%%%%%%%%%%%%%%%%%%%%%%%%%%%

%isso deve ser removido
\textit{$<$Retome brevemente o problema e o ODS definidos na Fase 1, indicando eventuais ajustes de escopo. Apresente o fluxo completo da solução, da entrada até a saída entregue, com uma figura da arquitetura. Justifique por que a abordagem multiagente é adequada ao problema.$>$}

%%%%%%%%%%%%%%%%%%%%%%%%%%%%%%%%%%%%%%%%
\section{Conjunto de Dados para a Avaliação}
\label{sec:dados}
%%%%%%%%%%%%%%%%%%%%%%%%%%%%%%%%%%%%%%%%

%isso deve ser removido
\textit{$<$Descreva a definição, a captura e a preparação dos conjuntos de dados: quais fontes são utilizadas, como são obtidas, qual o formato e quais restrições de uso se aplicam. Detalhe o pré-processamento e como os dados ficam disponíveis aos agentes em tempo de execução. Trate de questões de qualidade, viés e dados sensíveis, quando aplicável.$>$} %Explique como os dados serão coletados, as fontes utilizadas e os critérios para seleção. Detalhe o processo de preparação e limpeza dos dados, incluindo técnicas como normalização, remoção de duplicatas, e transformação de dados para adequação ao modelo.

%%%%%%%%%%%%%%%%%%%%%%%%%%%%%%%%%%%%%%%%
\section{Arquitetura Multiagente}
\label{sec:arquitetura}
%%%%%%%%%%%%%%%%%%%%%%%%%%%%%%%%%%%%%%%%

%isso deve ser removido
\textit{$<$Especifique cada agente do sistema: papel, entradas e saídas, escopo de decisão e critério de término da sua tarefa. Indique o modelo de linguagem atribuído a cada agente e justifique a escolha. Descreva as estratégias de raciocínio, planejamento e memória empregadas.$>$}

%%%%%%%%%%%%%%%%%%%%%%%%%%%%%%%%%%%%%%%%
\section{Orquestração e Uso de Ferramentas}
\label{sec:orquestracao}
%%%%%%%%%%%%%%%%%%%%%%%%%%%%%%%%%%%%%%%%

%isso deve ser removido
\textit{$<$Descreva como os agentes da Seção \ref{sec:arquitetura} atuam em conjunto: o padrão de orquestração adotado, o protocolo de comunicação, os limites de execução e o tratamento de falhas e divergências. Liste as ferramentas disponibilizadas aos agentes, a condição de uso de cada uma e como suas saídas são validadas. Indique os mecanismos de controle e os pontos de intervenção humana, quando houver.$>$}

%%%%%%%%%%%%%%%%%%%%%%%%%%%%%%%%%%%%%%%%
\section{Procedimentos de Avaliação}
\label{sec:avaliacao}
%%%%%%%%%%%%%%%%%%%%%%%%%%%%%%%%%%%%%%%%

%isso deve ser removido
\textit{$<$Defina os critérios de sucesso e as métricas adotadas, explicando como cada uma é calculada (por exemplo, precisão, correção do uso de ferramentas, custo e latência, avaliação por usuários alvo). Descreva o protocolo experimental: cenários de teste, condições controladas e \textit{baselines} de comparação. Explique como esses procedimentos garantem a robustez da metodologia e a validade dos resultados, indicando as limitações esperadas.$>$}

%%%%%%%%%%%%%%%%%%%%%%%%%%%%%%%%%%%%%%%%
\section*{Apêndice}
\label{sec:apendice}
%%%%%%%%%%%%%%%%%%%%%%%%%%%%%%%%%%%%%%%%

\textit{$<$Use se necessário. $>$}



%%%%%%%%%%%%%%%%%%%%%%%%%%%%%%%%%%
%          BIBLIOGRAFIA
%%%%%%%%%%%%%%%%%%%%%%%%%%%%%%%%%%
\newpage
\bibliographystyle{aasjournal}
\bibliography{references}{}



\end{document}
